%%=============================================================================
%% Inleiding
%%=============================================================================

\chapter{\IfLanguageName{dutch}{Inleiding}{Introduction}}%
\label{ch:inleiding}

De inleiding moet de lezer net genoeg informatie verschaffen om het onderwerp te begrijpen en in te zien waarom de onderzoeksvraag de moeite waard is om te onderzoeken. In de inleiding ga je literatuurverwijzingen beperken, zodat de tekst vlot leesbaar blijft. Je kan de inleiding verder onderverdelen in secties als dit de tekst verduidelijkt. Zaken die aan bod kunnen komen in de inleiding~\autocite{Pollefliet2011}:

\begin{itemize}
  \item context, achtergrond
  \item afbakenen van het onderwerp
  \item verantwoording van het onderwerp, methodologie
  
  
  \item onderzoeksdoelstelling
  
  \item onderzoeksvraag
  \item \ldots
\end{itemize}

Deze bachelorproef onderzoekt hoe een geautomatiseerd systeem voor security monitoring en incident response kan worden ontworpen om securitydata binnen SaaS-omgevingen effectief te centraliseren, analyseren en rapporteren in een DevOps-infrastructuur, en bouwt hier verder op aan de hand van een uitgewerkt proof-of-concept. Het onderzoek vertrekt van een casus bij Devinity, een SaaS-gericht bedrijf dat momenteel meerdere tools/platformen gebruikt voor security monitoring en incidentdetectie. Dit leidt tot problemen zoals vertraagde detectie, onvolledige rapporten en inconsistentie in het afhandelen van incidenten.  

De doelgroep bestaat uit het DevOps- en securityteam binnen Devinity, die baat hebben bij een gecentraliseerd en geautomatiseerd systeem dat de operationele efficiëntie verhoogt en risico’s beter beheersbaar maakt. Binnen de scope van dit onderzoek wordt gefocust op één representatief SaaS-platform en één DevOps-pipeline om de haalbaarheid en diepgang van het proof-of-concept te verzekeren.  

Bovendien benadrukken Gruhn \& Köhntopp (2021) dat security in DevOps-omgevingen een geïntegreerde aanpak vereist, waarbij monitoring en incident response nauw verweven zijn met de ontwikkeling- en deploymentprocessen, om snelle detectie en respons op beveiligingsincidenten te kunnen garanderen~\autocite{GruhnKohntopp2021}. 

\section{\IfLanguageName{dutch}{Probleemstelling}{Problem Statement}}%
\label{sec:probleemstelling}

Devinity BV krijgt dagelijks te maken met een concrete uitdaging binnen het vlak van security monitoring en incident response in de eigen SaaS- en DevOps-omgeving. Momenteel worden veel logs, events en alerts bekeken op verschillende afzonderlijke tools en platformen. 
Hierdoor vergroot dat een incident te laat wordt gedetecteerd, rapportages inconsistent of onvolledig zijn, en dat analyses manueel moeten uitgevoerd worden.

Voor de DevOps engineers en functional testers binnen Devinity betekent dit dat ze geen praktisch en gecentraliseerd overzicht hebben over de beveiligingsstatus van de infrastructuur. Daarnaast krijgt het managementteam geen duidelijke en uniforme rapportages, waardoor risicomanagement heel wat minder doelmatig verloopt.

Het ontbreken van een geautomatiseerde en gecentraliseerde oplossing heeft een directe impact op de analytische efficiëntie, de snelheid van incidentdetectie en de kwaliteit van de besluitvorming binnen Devinity.

Dit toont aan dat er binnen Devinity nood is aan een oplossing die logs, events en alerts uit de gebruikte platformen en omgevingen centraal verzamelt, deze analyseert en automatisch rapporteert. Op deze manier kan het bedrijft incidenten sneller detecteren, onnodige waarschuwingen beperken en een duidelijker overzicht krijgen van de beveiligingsstatus binnen het bedrijf.

\section{\IfLanguageName{dutch}{Onderzoeksvraag}{Research question}}%
\label{sec:onderzoeksvraag}

Hoe kan een geautomatiseerd systeem voor security monitoring en incident response worden ontworpen om securitydata binnen de SaaS-omgeving van Devinity effectief te centraliseren, analyseren en rapporteren in een DevOps-infrastructuur?

De \textbf{probleemstelling} wordt verder uitgewerkt aan de hand van volgende deelonderzoeksvragen:

\begin{enumerate}
    \item Welke uitdagingen ervaren SaaS-bedrijven zoals Devinity bij het gebruik van meerdere tools voor security monitoring en incidentdetectie?
    \item Hoe worden incidenten momenteel gedetecteerd en gerapporteerd, en welke tekortkomingen bestaan er in snelheid, correlatie en consistentie van rapporten?
    \item Welke securitydata (logs, events, alerts) worden beschikbaar gesteld door bestaande systemen en hoe worden deze momenteel gecentraliseerd en geanalyseerd?
    \item Welke metrics en KPIs worden nu gebruikt om de effectiviteit van security monitoring en incident response te meten, en welke zijn onvoldoende of inconsistent?
    \item Hoe beïnvloedt de huidige aanpak van monitoring en rapportering de operationele efficiëntie en het risicomanagement van SaaSbedrijven?
\end{enumerate}

De uitwerking van de \textbf{oplossing} gebeurt aan de hand van onderstaande deelonderzoeksvragen:

\begin{enumerate}
    \item Welke architectuur en technologieën zijn het meest geschikt om een geautomatiseerd securitymonitoringsysteem te ontwikkelen dat past binnen een DevOps-omgeving,
    \item Hoe kan eenvoudige anomaly-detectie op SaaS login-auditlogs, als aanvulling op rule-based detectieregels, bijdragen aan het identificeren van ongebruikelijke authenticatiepogingen?
    \item Welke data-integratieprocessen (APIs, agents, pipelines) zijn het meest efficiënt om logs en securitydata uit het geselecteerde SaaS-platform en de DevOps-pipeline te verzamelen en te normaliseren?
    \item Hoe kan automatische rapportgeneratie worden geïmplementeerd zodat rapporten zowel technische details als managementinzichten bevatten?
    \item Welke prestatie-indicatoren (KPIs) kunnen gebruikt worden om het succes van het systeem te meten, zoals detectiesnelheid, foutpositieven of responstijd?
    \item Hoe kan de beveiliging en betrouwbaarheid van het geautomatiseerde systeem zelf worden gegarandeerd binnen de bestaande DevOps-pipeline?
\end{enumerate}

\section{\IfLanguageName{dutch}{Onderzoeksdoelstelling}{Research objective}}%
\label{sec:onderzoeksdoelstelling}

Het doel van deze bachelorproef is het ontwikkelen van een proof-of-concept met het vermogen om securitydata uit een geselecteerde SaaS-platform en DevOps-omgeving automatisch te verzamelen en verwerken. Om ervoor te zorgen dat alle relevante gegevens centraal binnen het systeem terechtkomen kunnen verschillende technieken gebruikt worden, waaronder APIs, webhooks en logforwarders. Na het verzamelen van de data is het de bedoeling dat deze testimplementatie de gegevens structureert en normaliseert zodat ze op een efficiënte manier geanalyseerd kunnen worden.

Na het verwerken van deze gegevens zal er een automatisch rapport opgesteld worden dat de belanghebbenden uit het bedrijf voorziet van bruikbare informatie over de beveiligingsstatus om zo incidenten te voorkomen. Deze informatie kan eventueel ook KPIs bevatten om de effectiviteit van monitoring en incident response te meten.

Het systeem moet het dan ook mogelijk maken om sneller inzicht te krijgen in mogelijke incidenten, trends te herkennen in security-events en het overzicht van de infrastructuur te verbeteren, zodat het bedrijf geïnformeerde beslissingen kan nemen en de incidentrespons kan optimaliseren.

Concreet: het eindresultaat is een werkend prototype dat de voordelen van centralisatie, snelle detectie en consistente rapportage aantoont voor SaaS- en DevOps gerichte bedrijven zoals Devinity.

\section{\IfLanguageName{dutch}{Opzet van deze bachelorproef}{Structure of this bachelor thesis}}%
\label{sec:opzet-bachelorproef}

% Het is gebruikelijk aan het einde van de inleiding een overzicht te
% geven van de opbouw van de rest van de tekst. Deze sectie bevat al een aanzet
% die je kan aanvullen/aanpassen in functie van je eigen tekst.

De rest van deze bachelorproef is als volgt opgebouwd:

In Hoofdstuk~\ref{ch:stand-van-zaken} wordt een overzicht gegeven van de stand van zaken binnen het onderzoeksdomein, op basis van een literatuurstudie.

In Hoofdstuk~\ref{ch:methodologie} wordt de methodologie toegelicht en worden de gebruikte onderzoekstechnieken besproken om een antwoord te kunnen formuleren op de onderzoeksvragen.

% TODO: Vul hier aan voor je eigen hoofstukken, één of twee zinnen per hoofdstuk

In Hoofdstuk~\ref{ch:conclusie}, tenslotte, wordt de conclusie gegeven en een antwoord geformuleerd op de onderzoeksvragen. Daarbij wordt ook een aanzet gegeven voor toekomstig onderzoek binnen dit domein.